Дан целочисленный массив размера $N$. Вывести все содержащиеся в данном массиве нечетные числа в порядке возрастания их индексов, а также их количество $K$.
Дан целочисленный массив размера $N$. Вывести все содержащиеся в данном массиве четные числа в порядке убывания их индексов, а также их количество $K$.
Дан целочисленный массив размера $N$. Вывести вначале все содержащиеся в данном массиве четные числа в порядке возрастания их индексов, а затем --- все нечетные числа в порядке убывания их индексов.
Дан массив $A$ размера $N$ и целое число $K$ ($1 \leq K \leq N$). Вывести элементы массива с порядковыми номерами, кратными $K$: $A_K, A_{2 \cdot K}, A_{3 \cdot K}, \dots$ Условный оператор не использовать.
Дан массив $A$ размера $N$ ($N$ --- четное число). Вывести его элементы с четными номерами в порядке возрастания номеров: $A_2, A_4, A_6, \dots, A_N$. Условный оператор не использовать.
Дан массив $A$ размера $N$ ($N$ --- нечетное число). Вывести его элементы с нечетными номерами в порядке убывания номеров: $A_N, A_{N - 2}, A_{N - 4}, \dots, A_1$. Условный оператор не использовать.
Дан массив $A$ размера $N$. Вывести вначале его элементы с четными номерами (в порядке возрастания номеров), а затем --- элементы с нечетными номерами (также в порядке возрастания номеров): $A_2, A_4, A_6, \dots, A_1, A_3, A_5, \dots$ Условный оператор не использовать.
Дан массив $A$ размера $N$. Вывести вначале его элементы с нечетными номерами в порядке возрастания номеров, а затем --- элементы с четными номерами в порядке убывания номеров: $A_1, A_3, A_5, \dots, A_6, A_4, A_2$. Условный оператор не использовать.
Дан массив $A$ размера $N$. Вывести его элементы в следующем порядке: $A_1, A_N, A_2, A_{N - 1}, A_3, A_{N - 2}, \dots$
Дан массив размера $N$ и целые числа $K$ и $L$ ($1 \leq K \leq L \leq N$). Найти сумму элементов массива с номерами от $K$ до $L$ включительно.
Дан массив размера $N$ и целые числа $K$ и $L$ ($1 \leq K \leq L \leq N$). Найти среднее арифметическое элементов массива с номерами от $K$ до $L$ включительно.
Дан массив размера $N$ и целые числа $K$ и $L$ ($1 < K \leq L \leq N$). Найти сумму всех элементов массива, кроме элементов с номерами от $K$ до $L$ включительно.
Дан массив размера $N$ и целые числа $K$ и $L$ ($1 < K \leq L \leq N$). Найти среднее арифметическое всех элементов массива, кроме элементов с номерами от $K$ до $L$ включительно.
Дан массив $A$ размера $N$. Найти минимальный элемент из его элементов с четными номерами: $A_2, A_4, A_6, \dots$
Дан массив $A$ размера $N$. Найти максимальный элемент из его элементов с нечетными номерами: $A_1, A_3, A_5, \dots$
Дан массив из $N$ целых чисел. Выяснить, какое число встречается в массиве раньше --- положительное или отрицательное.
Дан массив из $N$ натуральных чисел. Создать массив из чётных чисел этого массива. Если таких чисел нет, то вывести сообщение об этом факте.
Дан целочисленный массив с количеством элементов $N$. Напечатать те его элементы, индексы которых являются степенями двойки ($1, 2, 4, 8, 16, \dots$).
Дан целочисленный массив размера $N$. Вывести все содержащиеся в данном массиве числа, кратные трём, в порядке возрастания их индексов, а также их количество $K$.
Дан целочисленный массив размера $N$. Вывести вначале все отрицательные числа в порядке возрастания их индексов, а затем --- все положительные числа в порядке убывания их индексов.
Дан массив $A$ размера $N$. Вывести его элементы в следующем порядке: $A_N, A_1, A_{N-1}, A_2, A_{N-2}, A_3, \dots$
Дан массив размера $N$ и целые числа $K$ и $L$ ($1 \leq K \leq L \leq N$). Найти произведение элементов массива с номерами от $K$ до $L$ включительно.
Дан массив размера $N$ и целые числа $K$ и $L$ ($1 < K \leq L < N$). Найти сумму элементов массива с номерами $1, \dots, K-1$ и $L+1, \dots, N$ (то есть сумму крайних элементов, исключая промежуток).
Дан массив $A$ размера $N$. Найти сумму минимального и максимального элементов из его элементов с четными номерами: $A_2, A_4, A_6, \dots$
Дан массив $A$ размера $N$. Найти номера (индексы) первого и последнего минимального элемента из его элементов с нечетными номерами: $A_1, A_3, A_5, \dots$
Дан целочисленный массив размера $N$. Вывести все элементы массива, которые больше предыдущего элемента (первый элемент не с чем сравнивать, поэтому не выводится).
Дан целочисленный массив размера $N$. Вывести все пары соседних элементов $(A_1, A_2), (A_2, A_3), \dots, (A_{N-1}, A_N)$, в которых оба числа четные.
Дан массив из $N$ целых чисел. Выяснить, сколько в массиве пар соседних элементов с разными знаками (положительное и отрицательное число). Нули не учитывать, считая, что знак нуля не определен.
Дан целочисленный массив с количеством элементов $N$. Напечатать те его элементы, индексы которых являются числами Фибоначчи ($1, 2, 3, 5, 8, 13, \dots$). Индексы, выходящие за пределы $N$, не обрабатывать.
Дан массив из $N$ целых чисел. Создать массив из чисел данного массива, которые оканчиваются на цифру 5. Если таких чисел нет, то вывести сообщение об их отсутствии.
